% v12 Threats, Discussion, Ethics, Conclusion.
\section{Threats to Validity}\label{sec:threats}
\textbf{Construct.} Every reference here scores location and change, never the concern: a correct start line does not show that a comment still applies, and a \lbl{changed} label does not show that a concern is affected. The concern-relative annotation takes the application of a suggestion as its reference for an addressed concern, which a suggestion can satisfy without settling the concern. Sampled anchors need not resemble the spans supervisors select; supervisors may favor short, repeated lines, which the stress suites show to be harder. Edits around the quoted range can matter to a concern and evade every quotation check (Section~\ref{sec:suggest}).

\textbf{Reference.} The real-history reference is a syntax-tree procedure, not semantic ground truth, and it excludes \rExcludedPct\% of candidates, mostly changed ones, so its accuracies are conditional on a decidable subset. Its weak step, same-type pairing, agreed with two blinded models on nearly every sampled target, its errors bound the deployed miss rate between \refdirBoundPct\% and \refcopyMissPctWorst\%, the upper end coming from copies it may take for moved code; its deletions are its least reliable decisions (Section~\ref{sec:refaudit}). The adjudicators are models, not people, and the prepared human packet is unannotated. The two references not labeled by us, the human-annotated benchmark (\cExtHumN{} regions) and the suggestion oracle (\wkUnits{} targets, externally authored but classified by our code), score file pairs, so moves between files are exercised only on our own reference.

\textbf{Internal.} The guards were designed from real-history failures, so their real-history results are retrospective (tier E1); AGi was specified after an earlier copy guard's student and compound results had been read (E2); the stress seed, the copy suite and the suggestion oracle are held out (E3, E4). Policy selection, the stratified analyses and the wrap mechanism read records that already existed. Protocols are local, timestamped records, not external preregistrations. CodeMapper and the language models are not fully deterministic (E5).

\textbf{Statistical and external.} Observations share files, commit pairs and projects; with nine projects, project-level intervals are coarse, and many secondary comparisons are unadjusted. The histories are Python; the suggestion oracle adds Go, Java, Rust and TypeScript, but as file pairs at a pull request's last head, not as resubmission chains. None of the populations consists of supervised student resubmissions. The language-model results are dated observations of free hosted models.

\section{Discussion}\label{sec:discussion}
\textbf{What was and was not shown.} The experiments show where anchored code is found, how its change is labeled and how often a request is withheld, on defined distributions; they are technical feasibility evidence. They do not show that supervisors re-review faster, revisit more of their earlier points or close fewer issues in error, or that students respond better. The direct test is a within-subjects study in which supervisors re-review the same revised submissions with a conventional interface and with \sys{}, measuring time to locate earlier comments, the share of earlier issues correctly revisited, erroneous closures and workload.

\textbf{Locator and layer.} Locating code is a matching problem on which stronger tools exist; deciding when a carried concern must be looked at again is a policy problem whose failures are properties of what a locator returns. That is why the same labels and guards lowered missed re-checks over every locator tested, and why the remaining failures (adjacent copies, copies added by the same revision, edits around an unchanged quotation, a deleted target placed on a look-alike) are the same for all of them. A model's concern judgment is a candidate annotation beside a request, untested with supervisors; no model output may close an issue.

\textbf{Configuring a deployment.} Four parts need no decision: resolution stays with the supervisor; requests are computed from the reviewed quotation, never from a locator's verdict; the reviewed quotation is kept beside the current one; and the strongest available locator is used, here CodeMapper. The copy guard is a choice: AGi avoids AG's requests on code shifted by insertions above it, while AG catches more copies added in the same revision that edits the original, so a course whose students copy and edit blocks together should prefer AG. Selecting thresholds or configurations per course from past data did not reliably beat a fixed CodeMapper+IG+AGi on held-out projects, and the identity floor trades wrong endpoints for lost ones without affecting missed changes. To make these choices from data, the portal should log, with consent, which carried issues a supervisor finds changed and which anchors a supervisor corrects.

\textbf{Ethics and data handling.} No student submission was used. The histories came from our projects and public repositories, the student-authored ones pseudonymized; source was parsed, never executed, and only public code was sent to hosted models. Portal telemetry requires opt-in consent that does not affect a student's access; a classroom study would need institutional approval.

\section{Conclusion}
Carrying review feedback safely is a different problem from relocating it. \sys{} keeps three properties apart, whether the reviewed code changed, whether the anchor still points at it and whether the concern is open, and states five requirements under which relocation evidence moves an issue only towards the supervisor; its layer applies the first three around any locator, under stated assumptions. The contribution is this contract and layer, not the matcher: CodeMapper locates code better than the portal's cascade on both references we did not label. Over every locator tested, the quotation labels and guards turned most hidden changes into requests at a measured cost in confirmations. What remains is specific: copies adjacent to the original or added by the same revision, edits that wrap an unchanged quotation, and requests that point at the wrong code. The implemented lifecycle never resolves an issue on its own but loses issues through deletion, dismissal keys and uncomposed versions; a stronger contract that closes these gaps is specified and model-checked, not yet implemented. The next steps are a consented study of supervisors on real resubmissions, human annotation of the prepared reference packet, and the stronger contract in the portal.

\section*{Acknowledgment}
\vAcknowledgment

% Generated by tools/v11/build_repro_tex.py; do not edit.
\section*{Reproducibility}
\begingroup\raggedright
\textbf{Pinned inputs.} Public real histories: Requests (\url{https://github.com/psf/requests}) at \texttt{3c7baaa0\allowbreak 5ab4aa49\allowbreak cc819ce3\allowbreak a819d260\allowbreak de19c157}; Flask (\url{https://github.com/pallets/flask}) at \texttt{d086db85\allowbreak 6be18725\allowbreak 5b8ec61e\allowbreak f4093573\allowbreak 93020f32}; Click (\url{https://github.com/pallets/click}) at \texttt{2247b35e\allowbreak a1c47c72\allowbreak 7d7a06e5\allowbreak 1fa280e1\allowbreak 2a863ff6}; HTTPX (\url{https://github.com/encode/httpx}) at \texttt{b5addb64\allowbreak f0161ff6\allowbreak bfe94c12\allowbreak 4ef76f6a\allowbreak 1fba5254}; Rich (\url{https://github.com/Textualize/rich}) at \texttt{9d8f9a37\allowbreak 2cc5916f\allowbreak d4781fec\allowbreak 207ced7d\allowbreak dac2f08f}, each fetched with Git depth 400, which follows every parent, and traversed by first parent (3,156 first-parent commits in all). Contextual tracker: \url{https://github.com/oselvar/lhdiff} at \texttt{5a9f7fd8\allowbreak 7193098c\allowbreak f90f7ee4\allowbreak 081f6d95\allowbreak f04d4cc7}. CodeMapper: \url{https://github.com/sola-st/CodeMapper} at \texttt{944fa91e\allowbreak 6b359322\allowbreak c1a91f25\allowbreak 5ed7f4df\allowbreak c556b68f}. GumTree: release v4.1.0 at \texttt{bd712afd\allowbreak a055f3b9\allowbreak 7dca0f33\allowbreak f2b73d36\allowbreak c712edfe}. Scale corpora: Django django-5.1/django-5.2 at \texttt{373cb303\allowbreak 7fe4e67a\allowbreak dbac9ac4\allowbreak 3340391e\allowbreak 859aa957}/\texttt{9e7cc2b6\allowbreak 28fe8fd3\allowbreak 895986af\allowbreak 9b7fc952\allowbreak 5034c1b0}; CPython cpython-v3.12.0/cpython-v3.13.0 at \texttt{0fb18b02\allowbreak c8ad5629\allowbreak 9d6a2910\allowbreak be0bab8a\allowbreak d601ef24}/\texttt{60403a54\allowbreak 09ff2c3f\allowbreak 3b07dd2c\allowbreak a91a7a3e\allowbreak 096839c7}. The 20 student-authored repositories are pseudonymized in this paper because they are student work; their URLs and full head hashes are in the selection manifest (SHA-256 \texttt{daa5e124\allowbreak 06698034\allowbreak 56ca3c69\allowbreak de8a1b43\allowbreak e8419ce3\allowbreak 54e9f775\allowbreak 0a3a4de8\allowbreak f7155767}). The four internal histories and the controlled corpus are private.

\textbf{Environment.} Python 3.12.14 (the compound-edit and multi-anchor stages ran under Python 3.11.4, recorded in their start files); NumPy 2.5.4 and Matplotlib 3.10.7 vendored for analysis and figures; CodeMapper and GumTree run in isolated environments whose dependency manifests and hashes are retained with the locator-comparison evidence. Timings come from one workstation (13th Gen Intel(R) Core(TM) i9-13900HX, Windows~11).

\textbf{Commands.} Each work package has a frozen protocol and a runner under \path{tools/v11/} (\path{run_wpb2.py}, \path{run_wph2.py}, \path{run_wpc.py}, \path{run_wpd_tests.py} and the real-history gate), which writes JSON with the SHA-256 of its code and inputs. \path{build_v11_numbers.py} turns those JSON files into every number in the text, \path{build_floats.py} and \path{build_effort_figure.py} build the remaining floats, and this section is generated by \path{build_repro_tex.py}. The macro file has SHA-256 \texttt{7b33f318\allowbreak 01bf8e0a\allowbreak b5a69119\allowbreak cb49ca3d\allowbreak 58837250\allowbreak 2325dd97\allowbreak 14f86ee7\allowbreak 919f1305} and records the hashes of its 27 input files. Artifact: \url{https://github.com/Noel-Alex/living-review-artifact}; its verifier rebuilds every macro from the shipped evidence. Internal histories and student repositories are not released.\par
\endgroup

